\section{总结与延伸}

\subsection{核心结论}

\begin{importantbox}{十二条可迁移结论}
\begin{enumerate}
  \item Hyperscaler primitives 与 application intent 之间仍需要编译层。
  \item 抽象必须创造新不变量，而不是只给控制台换皮。
  \item 产品对象的基数和生命周期决定资源映射是否经济。
  \item FDI 用框架语义生成 IR，再映射为 routing、cache 和 functions。
  \item ISR 按内容变化频率而非访问频率消耗 origin compute。
  \item Build vs. Buy 应把自研放在 insight 会复利的差异化层。
  \item Immutable deployment 与可变 pointer 让 preview、promote 和 rollback 可组合。
  \item Consumption pricing 只有可归因时才成为开发者反馈信号。
  \item Telemetry 要把版本、性能、错误和 spend 放在同一时间线。
  \item Soft cap 保护学习与连续性，hard cap 保护最坏支出但可能中断业务。
  \item AI workload 要区分 active CPU 与 I/O wait，compute density 不能靠固定 pod 数。
  \item Day 1 获得采用，Day 100 吸收复杂度，Day 1000 赢得运营信任。
\end{enumerate}
\end{importantbox}

\subsection{实践作业：设计一个 framework-aware AI deployment platform}

选择一个 AI web application，例如多模型聊天、coding agent、文档分析或生成式电商页面，完成以下设计：

\begin{enumerate}
  \item 定义 framework 能提供的应用语义，以及平台无法自动推断的业务约束；
  \item 设计 Build Output / IR，列出 assets、functions、routes、cache、tools 和 policies；
  \item 给出 immutable deployment、domain pointer、preview 和 rollback 模型；
  \item 将请求分为 CPU-bound、I/O-bound、streaming 和 background 四类，设计 capacity policy；
  \item 设计 telemetry schema，把 deployment version、latency、error、model token、CPU、memory 和 spend 关联起来；
  \item 设计 soft/hard controls，说明阈值、执行延迟、业务连续性和 abuse handling；
  \item 分别写出 Day 1、Day 100、Day 1000 的验收测试。
\end{enumerate}

\begin{knowledgebox}{验收标准}
高质量方案不会只写“用 serverless 自动扩缩容”。它应说明 framework intent 如何进入 IR，调度器如何区分 active compute 与 wait，成本如何归因到 operation，hard control 如何避免误伤，以及 rollback 遇到 stateful dependency 时如何处理。
\end{knowledgebox}

\subsection{拓展阅读}

\begin{enumerate}
  \item Vercel, \href{https://vercel.com/blog/framework-defined-infrastructure}{Framework-Defined Infrastructure}。
  \item Vercel, \href{https://vercel.com/blog/build-output-api}{Announcing the Build Output API} 与 \href{https://vercel.com/docs/build-output-api/v3}{Build Output API v3}。
  \item Vercel, \href{https://vercel.com/blog/life-of-a-request-application-aware-routing}{Life of a Vercel request: Application-aware routing}。
  \item Next.js, \href{https://nextjs.org/docs/app/guides/incremental-static-regeneration}{Incremental Static Regeneration}。
  \item Vercel, \href{https://vercel.com/blog/introducing-fluid-compute}{Introducing Fluid Compute} 与 \href{https://vercel.com/docs/fluid-compute}{current documentation}。
  \item Vercel, \href{https://vercel.com/docs/spend-management}{Spend Management} 与 \href{https://vercel.com/docs/observability}{Observability}。
  \item v0, \href{https://v0.app/docs}{Documentation}。
\end{enumerate}

\subsection{建议阅读路线}

先读 FDI 与 Build Output API，画出“源码—框架语义—IR—部署”的编译边界；再读 request routing 与 ISR，追踪一个请求怎样命中 metadata、cache、function 和 origin。随后阅读 Spend Management 与 Observability，把监控项映射到课堂的 consumption feedback。最后阅读 Fluid Compute，并用 active CPU / wall time 比例解释为什么 AI request 改变了 serverless 的调度目标。阅读每个官方页面时，统一记录\textbf{输入契约、生成或管理的资源、可观测信号、失败/回滚路径，以及当前版本与历史课堂的差异}；若页面没有明确 failure mode，再主动补写可能的控制面、数据面和计量风险。

\end{document}
